% ============================================================
% f2014_shandong_science_q07_histogram.tex
% 2014 山东卷（理科）第 7 题：舒张压数据的频率分布直方图
% 来源 PDF：普通高考/2014/2014山东理.pdf
% 原页：第 1 页（A3 横排）左栏，印面 907，第 7 题题干下方
% 各矩形高（频率/组距）：0.24、0.16、0.36、0.16、0.08，和为 1。
% ============================================================
\begin{tikzpicture}[>=Stealth, line width=0.7pt, font=\small,
  x=1cm, y=1cm]
  % 坐标轴
  \draw[->] (0.60,0) -- (7.35,0) node[below right] {舒张压/kPa};
  \draw[->] (0.60,-0.05) -- (0.60,4.05);
  \node[below left] at (0.60,0) {$0$};
  % 纵轴标题：频率 / 组距（上下两行，原图无分数线）
  \node at (0.62,4.62) {频率};
  \node at (0.62,4.16) {组距};
  % x 轴断开符号
  \draw (0.86,0) -- (0.98,0.10) -- (1.10,-0.10) -- (1.22,0.10) -- (1.34,0);
  % 直方图：组距 1，映射 12→1.90，13→2.90，…，17→6.90
  \draw (1.90,0) rectangle (2.90,2.40);
  \draw (2.90,0) rectangle (3.90,1.60);
  \draw (3.90,0) rectangle (4.90,3.60);
  \draw (4.90,0) rectangle (5.90,1.60);
  \draw (5.90,0) rectangle (6.90,0.80);
  % 纵轴刻度与虚线
  \foreach \h/\y in {0.08/0.80, 0.16/1.60, 0.24/2.40, 0.36/3.60}
  {
    \draw (0.54,\y) -- (0.66,\y);
    \node[left] at (0.52,\y) {$\h$};
    \draw[dashed] (0.66,\y) -- (6.90,\y);
  }
  % 横轴刻度
  \foreach \x/\l in {1.90/12, 2.90/13, 3.90/14, 4.90/15, 5.90/16, 6.90/17}
  {
    \draw (\x,0.05) -- (\x,-0.05);
    \node[below] at (\x,-0.02) {$\l$};
  }
\end{tikzpicture}
